\textbf{Chapter 1 -- Introduction}

\textbf{1.1 Background and Motivation}

A private investor holding an S\&P 500 asset faces a similar decision
each trading day: buy more, sell part of the position, or do nothing. On
most days, the choice may have little effect on the overall outcome.
During a heavy market move, however, the decision can have a much larger
impact.

March 2025 is one of such periods studied in this dissertation. It was
the worst month for the buy-and-hold strategy, which lost \$398.41 from
an initial \$10,000 investment, approximately 4\% of the starting
capital. An investor who reduced their position before the decline could
have avoided part of this loss. However, reducing positions too often
would also mean missing gains during the other months, which produced an
average monthly gain of more than \$200.

The problem is therefore not just when to invest, or deciding when the
market will rise or fall, but how to make a series of decisions about an
asset's position as market conditions change, as each decision made
affects the position long-term and the consequences of that action will
only become clear several days or weeks later.

This makes trading a sequential decision-making problem, which is one
reason reinforcement learning (RL) is appropriate for the task. RL
allows an agent to learn by interacting with an environment, observing
the results of its actions, and receiving rewards over time based on
those results. In essence, it learns a policy/strategy that makes better
decisions to receive more rewards.

Trading has several characteristics that make it suitable for this
machine learning approach. The agent observes information about the
market and its portfolio, selects from a set of possible actions such as
buying, selling, or holding, and receives a measurable reward based on
the change in portfolio wealth. As a result, reinforcement learning has
been increasingly applied to financial trading.

However, another important problem arises. The Stock/Equity markets
generally rise over the years rather than fall. Meaning that, in a
rising market, a strategy that simply remains invested without any
pulling out can generate a strong return without using any market
information. A learned agent that produces a positive return therefore
has not necessarily demonstrated, by itself, that it has learned a
useful trading strategy in such markets.

Two questions become relevant. First, can the learned agent outperform
the simple buy-and-hold strategy with which it is being compared?
Second, does the agent actually use the information available to it when
deciding what action to take? These questions are important when
evaluating on data not used during training.

This dissertation therefore focuses on both the financial performance of
reinforcement learning agents and whether their decisions show evidence
of using the information provided to them.

\textbf{1.2 Problem Statement}

The main problem this dissertation addresses is whether reinforcement
learning algorithms can learn a trading policy for a portfolio ---
containing only one asset in this case --- that beats a simple benchmark
strategy on unseen data.

However, financial performance alone is not enough to determine that an
agent has actually learned a useful trading strategy. For example, an
agent that buys whenever it can, would perform well in a generally
rising market without using any of the market information provided to
it. A positive return from such a policy could be incorrectly
interpreted as proof that the agent learned when to buy, sell, or remain
out of the market, whereas that was not the case.

Evaluating a learned trading agent therefore requires a second
perspective alongside financial performance: \textbf{behaviour}. The
study must analyse whether the agent\textquotesingle s actions are a
result of changes in the observed market and portfolio state. This helps
to differentiate a policy that genuinely responds to the information it
receives from one that simply follows a fixed trading pattern.

This dissertation therefore addresses two interrelated problems: whether
reinforcement learning can produce better financial performance than
simple benchmarks, and whether that performance stems from a learned
policy that actually uses the information available to it.

\textbf{1.3 Research Questions and Objectives}

The aim of this dissertation is to evaluate three reinforcement learning
algorithms --- REINFORCE, Deep Q-learning (DQN), and Proximal Policy
Optimization (PPO) --- and determine whether they can learn profitable
and genuinely intelligent trading strategies for the SPY exchange-traded
fund.

The study is structured around four research questions:

\begin{enumerate}
\def\labelenumi{\arabic{enumi}.}
\item
  \textbf{Performance:} Do the learned agents outperform simple
  benchmarks on unseen data after transaction costs are taken into
  account?
\item
  \textbf{Behaviour:} Do the learned agents use the market and portfolio
  information provided in their state representation when making trading
  decisions?
\item
  \textbf{State representation:} Does adding a drawdown feature to the
  main nine-feature state change the behaviour learned or the financial
  performance?
\item
  \textbf{Reward:} Does a risk-aware reward improve the balance between
  return and drawdown, rather than simply encouraging the agent to trade
  less?
\end{enumerate}

These research questions are touched upon through seven measurable
objectives, each traceable to a specific part of the dissertation

\textbf{O1.} Build a reproducible daily trading environment for SPY
using monthly episodes, transaction fees, and action masking, and verify
the environment (Chapter 4).

\textbf{O2.} Implement REINFORCE, DQN, and PPO and verify their
implementations before running the main trading experiments (Chapters 3
and 4).

\textbf{O3.} Calibrate the fixed trade size on validation data so the
agents' exposure to the market is not limited, providing a fair
comparison to the buy-and-hold benchmark (Section 5.3).

\textbf{O4.} Apply behavioural analysis to determine whether the learned
policies change their actions in response to the observed state (Section
5.4).

\textbf{O5.} Measure the effect of adding drawdown to the state while
keeping the other experimental conditions fixed (Section 5.5).

\textbf{O6.} Compare the final learned agents with the benchmark
strategies over 24 held-out test months, using six independently trained
random seeds for each configuration (Section 5.6).

\textbf{O7.} Examine whether results change when a different reward
formula is used (Section 5.7).

These objectives create a structured way to evaluate both the financial
performance and the learned behaviour of the three reinforcement
learning methods.

\textbf{1.4 Scope}

The study is deliberately narrowed to one asset to make the experiments
more controlled and easier to analyse.

All experiments use the SPY asset, an exchange-traded fund that tracks
the S\&P 500 index, with daily data from 2018 to 2025. The data are
divided chronologically into training (2018--2022), validation (2023),
and held-out for testing periods (2024--2025).

Each monthly episode starts with \$10,000 of capital and no asset
holdings, and each transaction to open a position incurs a 0.05\% fee.
The available actions are discrete: the agent can buy a fixed fraction
of the initial capital, sell the same amount, or hold. The action
masking prevents invalid actions, such as attempting to sell when no
position is held.

Three reinforcement learning algorithms are evaluated. REINFORCE and
Deep Q-learning are implemented from scratch, while PPO is implemented
using a popular library.

Several areas are intentionally outside the scope of this study. These
include multi-asset portfolios, continuous position sizing, intraday
trading data, and live deployment. Keeping the study within these
boundaries limits the number of factors that can affect the results,
making it easier to identify which specific factor causes a change in
performance or behaviour during testing rather than multiple changes at
once.

\textbf{1.5 Contributions}

This dissertation makes four main contributions:

\textbf{1. A framework for evaluating trading agents.}

The study uses a controlled process in which experiment variables, such
as trade size and transaction cost, are set using validation data and
finalized before moving on to the held-out test data. The trading
environment, state features, rewards, and benchmark strategies are also
verified through tests to identify and correct problems before final
evaluation,

\textbf{2. A behavioural analysis for testing state utilization.}

A separate analysis was developed to determine whether a policy learned
actually responds to the information in its state, rather than simply
following a fixed pattern of available actions.

In this evaluation, steps are grouped by their action mask. When a
policy selects different actions at steps with the same mask pattern,
the difference cannot be explained by the available actions alone,
thereby demonstrating a change in its action based on the current
observed state. The analysis does not identify which individual state
feature drives this dependence, nor does it make claims about the
learned network\textquotesingle s internal representations.

This reveals a difference between algorithms that would not be seen from
financial return figures alone: all six Deep Q-learning seeds were
state-dependent, while none of the twelve policy-gradient seeds
(REINFORCE and PPO) were (Section 5.4).

\textbf{3. An explanation for the negative results.}

None of the three learning algorithms outperformed buy-and-hold across
the 24 held-out test months. The behavioural analysis helps explain this
outcome. The policies that achieved returns closest to the benchmark,
REINFORCE, generally did so through a fixed pattern of high exposure,
while the algorithm that did use its state information, Deep Q-learning,
reduced market participation rather than successfully understanding
market conditions to perform better.

Although the overall result is negative, evidence from the six training
seeds per configuration, behavioural analysis, and peripheral
experiments on state representation, reward function, and transaction
costs supports it as a finding rather than a failure.

\textbf{1.6 Dissertation Structure}

The dissertation is organized as follows. Chapter 2 reviews existing
research on reinforcement learning for financial trading and portfolios
and explains how this study's evaluation approach relates to previous
work.

Chapter 3 presents the mathematical formulation of the portfolio trading
problem, describing the trading environment as a Markov decision
process; the state, action, and reward formulas; the three reinforcement
learning algorithms; and their optimisation functions.

Chapter 4 describes the system implementation, data pipeline,
verification process, trading environment, and the corrections made
during development.

Chapter 5 presents the experiment results and analyses them in relation
to the four research questions. It reports the trade-size calibration,
behavioural analysis, benchmark comparison, state-representation and
reward experiment, and transaction-cost calibration.

Finally, Chapter 6 concludes the dissertation by summarising the main
findings, critically evaluating the work, discussing its limitations and
possible directions for future research.

\textbf{\hfill\break
}

\textbf{6.4.2 Monthly Episode Structure}

The trading environment treated each calendar month as a separate
episode. Every episode started with the same initial capital of
\$10,000, and any remaining position was closed at the end of the month.
As a result, the results do not represent a single portfolio that
continuously compounds over the full evaluation period as a real
portfolio would.

This design made the experiments easier to control by keeping the
episode length and starting conditions consistent. However, it also
means that gains or losses from one month did not affect the starting
capital or portfolio position in the following month, which does not
represent a real-world portfolio. A longer continuous portfolio
environment could therefore lead to a different behaviour, especially
for policies that require accumulated wealth or changes in portfolio
value over time.

\textbf{5.2 Experiment Protocols and Evaluation Criteria}

This section describes how the experiments in this chapter were
conducted and how the reported results should be interpreted.

\textbf{5.2.1 Experimental Configurations}

Each configuration trained REINFORCE, Deep Q-learning (DQN), and
Proximal Policy Optimisation (PPO) separately for 80,000 steps per seed
in the portfolio environment,

To account for randomness in reinforcement learning, six random seeds
were used, with an initial capital of \(C_{0} = \$ 10,000\) at the start
of every monthly episode. This means six independent training runs all
with the same configuration, allowing us to analyse the consistency of
each algorithm's results rather than relying on a single training run.

During evaluation, all policies acted deterministically: policy-gradient
algorithms select the action with the highest probability, while
value-based algorithms select the action with the highest estimated
value.

The 2023 validation period was used to make every experimental decision
reported in this chapter. This included selecting the trade size in
Section 5.3, the risk-penalty weight in Section 5.7, and the transaction
fee in Section 5.8. These choices were fixed before using them for the
main experiments in the 24 held-out test months from 2024--2025.

There are three main learned experiments in this chapter:

\begin{itemize}
\item
  The primary experiment evaluating the nine-feature state with the
  wealth-change reward.
\item
  The ten-feature state that adds drawdown while keeping the same
  wealth-change reward.
\item
  The third experiment using the ten-feature state with the risk-aware
  reward, which actually requires drawdown.
\end{itemize}

\textbf{5.2.2 Evaluation Metrics}

There are five financial measures used in each experiment report. They
are averaged over the episodes and, for learned agents, across the six
training seeds:

\begin{itemize}
\item
  Mean wealth change \((\Delta W)\) reported in dollars per month;
\item
  Mean exposure \((\Delta v_{t})\), which represents the proportion of
  initial capital invested in the asset and is averaged across trading
  days;
\item
  Trades per month;
\item
  Mean episode drawdown \((\Delta{dd}_{t})\);
\item
  and the monthly Sharpe ratio, a financial metric that calculates the
  average return relative to the variation of those daily returns during
  the month.
\end{itemize}

The tables also report the \textbf{seed spread}, which is the difference
between the best and worst result across the six seeds. This is included
because there might be large differences between training runs which the
average alone may hide. A large spread indicates that the algorithm did
not produce a consistent trading behaviour across the different runs.

Financial performance is then followed by a behavioural analysis, as
returns alone cannot show whether a policy actually used the information
in its state. For example, a policy that buys simply because buying is
allowed can perform well during a rising market without responding to
any market information.

This led to \textbf{state-dependence analysis,} which was used to test
whether a policy\textquotesingle s actions changed when its observed
state changed while the valid actions remained the same. For instance,
if a policy selected different actions at different steps with the same
actions available, the difference could be attributed to the observed
state. The tables report how many of the six seeds met this condition.

\textbf{5.2.3 Benchmark Strategies}

Two non-learning strategies are used as benchmarks. They use the same
trading environment, portfolio conditions, and transaction fees as the
learned agents.

\textbf{Buy-and-hold} invests the full initial capital on the first
trading day and holds the position until the end of the month, when the
position is closed. This is the main benchmark as it is a simple passive
investment strategy.

\textbf{Always-buy} \(\mathbf{(0.25}\mathbf{C}_{\mathbf{0}}\mathbf{)}\)
purchases a fixed trade amount equal to 25\% of the initial capital
whenever buying is available, continuing until it has fully invested in
that asset. This strategy is useful for comparing against policies that
achieve their returns through repeated buying rather than by using
information from its observed state, as it also does not use any market
or portfolio information when deciding whether to buy.

\textbf{4.4.3 Benchmark Implementation}

The learned agents are compared with benchmark strategies to provide
clear grounds for assessing their performance. These benchmarks are
simple trading behaviours that do not require reinforcement learning.

The main benchmark is a buy-and-hold strategy where the initial capital
is used to purchase the asset at the start of the episode. The position
is then held until the final trading step, when it is closed. The same
transaction fee and position closure used for the learned agents apply
to this benchmark.

This approach is different from gradually purchasing the asset using the
same fixed trade amount available to the learning agents. A gradual
purchasing strategy may take several trading days to build a position
and, if the fixed trade size is small, it may not become fully invested
before the monthly episode ends.

For this reason, the gradual purchase strategy is treated as a separate
reference strategy rather than as buy-and-hold. The conventional
buy-and-hold strategy gives the learned agents a straight performance
target.

The benchmark strategies use the same portfolio conditions, transaction
costs, and evaluation procedures as the learned agents. This keeps the
comparison consistent and reduces the chance that differences in results
stem from inaccuracies in calculating portfolio performance rather than
from the strategies themselves.

\textbf{\hfill\break
}

\textbf{Statement of Originality}

I confirm that this report and the work it presents are my own. Any
ideas, data, images or text taken from the work of others, whether
published or unpublished, have been clearly identified and appropriately
acknowledged in the text and bibliography.

I confirm that this report has not been submitted, either in whole or in
part, for any other academic degree or professional qualification.

I understand and agree that the University may submit this work to
TurnitinUK to check its originality.

\textbf{Acknowledgements}

First and foremost, I would like to thank God for giving me the idea for
this project, the strength, wisdom and perseverance to see it through.
What started as an initial idea developed significantly throughout this
journey, and I am grateful to everyone who played a part in bringing it
to where it is today. It is only by His Grace to Start, Continue, and
Finish that this was possible

I would like to thank the professors who helped me narrow down the
original idea to shape it into a research problem that could be properly
investigated. I am especially grateful to my supervisor, Dr. Nguyen, for
his guidance throughout the project. His feedback helped me streamline
the work, question my assumptions, and keep the research focused on what
evidence could support. More importantly, his guidance helped me develop
a much stronger understanding of reinforcement learning and how it can
be applied to a problem as challenging as financial trading.

I would also like to express my deepest gratitude to my family,
especially my mum and dad, for their constant encouragement and support.
They kept pushing me to keep going until the final stage. I am equally
grateful to my uncle, aunties, grandma, church circle, and everyone else
who took the time to check in on me throughout the journey. Their
encouragement and help, even in small moments, but especially in sick,
tired moments, all meant a great deal to me.

Finally, I would like to give a special shoutout to my friends from the
cohort: Katherine, Joel, Yumna, Akeeb, David, and Dr. Tim. The laughs,
outings, and support throughout the year are unforgettable. Having
people around who were going through the same experience made the long
parts of the journey so much fun, shorter, and easier to get through.

To everyone who contributed to this journey in one way or another, thank
you. I am genuinely grateful.

\textbf{Abstract}

An investor holding a stock-index fund faces the same decision each day:
buy more, sell part of the position, or do nothing. This dissertation
asks whether reinforcement learning, which learns by using rewards via
trial and error, can make these decisions more effectively than the
simple strategy of buying and holding permanently.

Three algorithms are examined: REINFORCE and Deep Q-learning (DQN),
implemented from scratch, and Proximal Policy Optimisation (PPO), from
an open-source library. Each agent trades the S\&P 500 exchange-traded
fund (SPY) in a simulated environment using monthly episodes based on
daily data from 2018 to 2025. The data is split chronologically into
training data from 2018--2022, validation data from 2023, and held-out
test months from 2024--2025.

The environment includes the trade size, transaction fees, and action
validity. Experimental decisions, such as using a trade size equal to
one quarter of the initial \$10,000 capital, were selected using the
validation data before being applied to the test data. Each
configuration was trained using six different random seeds.

None of the learned agents beat buy-and-hold on the held-out test
period. REINFORCE earned the highest mean return at +\$171.51 per month,
compared with +\$180.25 for buy-and-hold. PPO earned +\$127.35, while
DQN earned +\$74.68.

Behavioural analysis provides an important explanation for these
results. All six DQN seeds changed their actions in response to the
observed state, whereas none of the twelve policy-gradient seeds
(REINFORCE and PPO) did so. The returns earned by these policy-gradient
algorithms were therefore largely tied to fixed buying patterns that
benefited from the generally rising test period, rather than clear
evidence of depending on market information to make informed decisions.

Two further experiments supported this interpretation. Adding drawdown
to the state produced little change in either behaviour or performance
for any of the algorithms. A risk-aware reward reduced
DQN\textquotesingle s drawdown by 41\%, but only by cutting earnings by
32\% and approximately halving market exposure.

The main contribution of the study is therefore an evaluative framework
that uses calibrated trade sizes, simple benchmarks, behavioural
analysis, and verified implementations, creating a reliable basis for
interpreting both successful and negative results from reinforcement
learning in trading.
